\section{The sharp residue receiver for fixed arithmetic heat tests}
\label{sec:residue-fixed-receiver}

The supplied \emph{Regular-inverse contact defects and their arithmetic
return}, Theorem \texttt{thm:residuejet}, proves the three-coefficient
receiver for collision energy and exhibits the missing coefficient at a
double collision. The result here carries that coefficient into every
fixed-test layer HR4 and the local determinant correction HR14 of
\texttt{FIXED\_ARITHMETIC\_HEAT\_FLAGS.tex}. The full flag source
\texttt{HEAT\_TRACE\_FLAG\_RESPONSE.tex}, HF1--16, is also retained.
Those two complete local sources and the supplied theorem were read.
The heat convention is Brad Rodgers and Terence Tao,
\href{https://arxiv.org/abs/1801.05914v5}{\emph{The de Bruijn--Newman
constant is non-negative}}, author equations \texttt{hoz},
\texttt{htdef}, and \texttt{sas}, as identified in HR and HF. This
component does not claim an independent reading of their entire paper.
All receiver maps, thresholds, and finite heat examples used below are
proved here.

\paragraph{RJ1. Original coordinates and marked data.}
Let $H(t,x_0+x)$ be a jointly holomorphic solution of
$\partial_tH=-\partial_x^2H$, of real type for real $t,x$, with an
actual zero of order $m\ge1$ at $(0,x_0)$. Retain
\[
 H(0,x_0+x)=x^m u(x),\qquad
 u(x)=u_0+u_1x+u_2x^2+\cdots,\quad u_0\ne0,
 \qquad \rho=\frac12+\frac{ix_0}{2},\quad s=\rho+\frac{ix}{2}.
 \tag{RJ1}
\]
In particular the actual arithmetic family may be used, with
$H_0(z)=\xi_{\mathrm R}(1/2+iz/2)/8$. No multiplicity of an
actual arithmetic zero is assigned by the following construction.
Write
\[
 c=\frac{u_1}{u_0},\qquad d=\frac{u_2}{u_0},\qquad
 b=2d-c^2=(\log u)''(0).
 \tag{RJ2}
\]
Thus $c,d,b$ are exactly the HR1 coefficients. The letter $b$ in
RJ2 is not the letter $b$ used for the coefficient $u_2/u_0$ in the
supplied energy paper: that paper's $a,b,\kappa$ are respectively
$c,d,b$ here. We keep the entire $u$, not a substituted polynomial
unit, when defining the receiver.

For $n\ge1$ define a coordinate-marked algebra and functional by
\[
 A_n=\mathbb C[x]/(x^n),\qquad
 \lambda_n(f)=\operatorname{Res}_{x=0}
          \frac{f(x)}{x^n u(x)}\,dx,\qquad
 B_n(f,g)=\lambda_n(fg).
 \tag{RJ3}
\]
Here ``marked'' specifies the class of the original $x$, the identity,
and the given functional; arbitrary changes of coordinate are not
identified with the same observation. The functional is independent
of the polynomial representative: adding $x^nh$ changes its integrand
by the holomorphic function $h/u$. The actual collision receiver is
$A_m$. Other values of $n$ record analytic thickenings or quotients,
not a changed zero multiplicity or a changed heat clock.

\paragraph{RJ2. What the entire residue pairing records.}
Put $u(x)^{-1}=\sum_{k\ge0}q_kx^k$. Direct extraction of the
coefficient of $x^{-1}$ gives
\[
 \lambda_n(x^a)=q_{n-1-a}\quad(0\le a<n),\qquad
 B_n(x^a,x^b)=
 \begin{cases}q_{n-1-a-b},&a+b<n,\\0,&a+b\ge n.\end{cases}
 \tag{RJ4}
\]
Reversing the columns of this matrix makes it triangular with diagonal
$q_0=u_0^{-1}$. Therefore
\[
 \det B_n=(-1)^{n(n-1)/2}u_0^{-n}\ne0.
 \tag{RJ5}
\]
This proves perfection with its full matrix, including all coefficients
of $u$ that can enter at this order. Conversely, $B_n(1,f)=\lambda_n(f)$,
so knowing the entire marked pairing gives exactly $q_0,\ldots,q_{n-1}$.
The identities $u(x)\sum q_kx^k=1$ recover
\[
 u_0=q_0^{-1},\qquad
 u_k=-q_0^{-1}\sum_{i=1}^k q_i u_{k-i}\quad(1\le k<n).
 \tag{RJ6}
\]
Hence two units give the same marked receiver at order $n$ if and
only if their jets agree modulo $x^n$. This is an exact equivalence;
perfection alone does not add any unobserved coefficient.

\paragraph{RJ3. Every comparison map in the residue tower.}
For $N=n+r$, $r\ge1$, let $\pi_{N,n}:A_N\to A_n$ be the quotient
algebra map, and define
\[
 \jmath_{n,N}(f)=x^r\widetilde f\pmod{x^N},
 \qquad
 0\longrightarrow (x^n)/(x^N)\longrightarrow A_N
       \xrightarrow{\pi_{N,n}}A_n\longrightarrow0.
 \tag{RJ7}
\]
A different lift changes $x^r\widetilde f$ by a multiple of $x^N$,
so $\jmath$ is well defined. Its image is $(x^r)$ and its kernel is
zero, because the coefficients of $x^r,\ldots,x^{N-1}$ are precisely
the original $n$ coefficients. The kernel of $\pi$ is $(x^n)$,
of dimension $r$, identified with $A_r$ as a module by multiplication
by $x^n$. Give $A_n$ the $A_N$ action through $\pi$. Then
$\jmath$ is an $A_N$-module map, since
$a x^r\widetilde f=x^r\widetilde{\pi(a)f}$ modulo $x^N$.
For $r>0$ it is not a unital algebra map: it sends $1$ to $x^r$.

The full adjunction identity is
\[
 \boxed{\lambda_n(f)=\lambda_N(\jmath_{n,N}f),\qquad
 B_n(\pi_{N,n}a,f)=B_N(a,\jmath_{n,N}f).}
 \tag{RJ8}
\]
Both statements follow from cancellation of $x^r$ in the residue
denominator; the second also follows from the module identity. Thus
under the perfect dualities $B_N,B_n$, the transpose of the quotient
map is exactly $\jmath$. In particular
\[
 (\ker\pi_{N,n})^{\perp_{B_N}}=\operatorname{im}\jmath_{n,N}=(x^r).
 \tag{RJ9}
\]
The inclusion follows from RJ8. Both sides have dimension $n$ by
perfection and the already computed kernel dimension, proving equality.
Composition retains the whole tower:
$\pi_{N,n}\pi_{P,N}=\pi_{P,n}$ and
$\jmath_{N,P}\jmath_{n,N}=\jmath_{n,P}$ for $P\ge N\ge n$.
These formulas are literal quotient composition and multiplication
by $x^{P-N}x^{N-n}$.

All support labels can be retained at these maps. On the original
full-support carrier $\{(0,\ell):\ell\in L\}\cup
\{(a,1_L):a\in A_n\}$, a linear map $\phi$ above sends
$(a,1_L)$ to $(\phi(a),1_L)$ and fixes every $(0,\ell)$.
This is well defined on the common top-supported zero. It preserves
addition because amplitudes add and labels join. Scalar action is
preserved with the module actions specified above. A nonzero element
in the kernel of $\pi$ becomes $(0,1_L)$, retaining its top label;
it does not become external absence. The quotient is also
multiplicative, since $\pi$ is an algebra map and labels meet.
Only the asserted module structure is preserved by $\jmath$.
The pairings in RJ8 retain the same output support on both sides:
the input labels have not changed, including when the scalar residue
is zero.

\paragraph{RJ4. The least residue orders for the original unit.}
For every $n\ge3$ write
\[
 L_0=\lambda_n(x^{n-1}),\quad L_1=\lambda_n(x^{n-2}),\quad
 L_2=\lambda_n(x^{n-3}),\qquad A=\frac{L_1}{L_0},\quad
 B=\frac{L_2}{L_0}.
 \tag{RJ10}
\]
The same $L_0,L_1$ exist at $n=2$. Equations RJ4--6 give
\[
 L_0=u_0^{-1},\qquad L_1=-cu_0^{-1},\qquad
 L_2=(c^2-d)u_0^{-1},
 \qquad\boxed{c=-A,\quad d=A^2-B,\quad b=A^2-2B.}
 \tag{RJ11}
\]
The tower proves these are independent of the chosen receiver order
once that order is large enough. Order $2$ is the least order that
determines $c$ for all units: $1+\alpha x$ has fixed order-$1$
receiver and arbitrary $c=\alpha$. Order $3$ is the least order
that determines either $d$ or $b$: $1+c x+d x^2$ has the same
order-$2$ receiver as $d$ varies, but both $d$ and $b=2d-c^2$
vary. These are comparisons of whole functionals and whole perfect
pairings, not selected entries alone. RJ9 describes the exact
duality map between these objects; no inference from a quotient's
nonidentity is used to deny their relation.

\paragraph{RJ5. Heat expansion in the original time.}
We give the analytic calculation needed to apply this receiver,
including the constants in HF and HR. Let
\[
 P_k(y)=\sum_{a=0}^{\lfloor k/2\rfloor}
       \frac{(-1)^ak!}{a!(k-2a)!}y^{k-2a}.
 \tag{RJ12}
\]
Coefficient differentiation gives
$P_k'=kP_{k-1}$,
$P_{k+1}=yP_k-2kP_{k-1}$, and
$2P_k''-yP_k'+kP_k=0$.
All roots of $P_k$ are real and simple. One proof uses
$P_k=(-2)^ke^{y^2/4}(d/dy)^ke^{-y^2/4}$, obtained by coefficient
differentiation. Integration by parts $k$ times makes $P_k$
orthogonal to every polynomial of degree less than $k$ against
$e^{-y^2/4}dy$; every boundary term vanishes by Gaussian decay.
If $P_k$ had fewer than $k$ sign-changing real roots, multiply the
linear factors at those roots. The resulting polynomial has degree
less than $k$, and its product with $P_k$ has one sign, is not zero,
and is integrable against that strictly positive weight. Its integral
cannot vanish, a contradiction. Thus there are exactly $k$ distinct
real roots. Denote those of $P_m$ by $r_1,\ldots,r_m$, and put
$p_a=\sum_\nu r_\nu^a$, including $p_0=m$.

Set $t=\epsilon^2$, $x=\epsilon y$ for the root calculation only.
The PDE forces
$\partial_t^a\partial_x^kH(0,x_0)=(-1)^a
\partial_x^{k+2a}H(0,x_0)$.
In the convergent two-variable Taylor series, every weighted degree
$k+2a<m$ therefore vanishes. Grouping the finitely many terms at each
of the next three weighted degrees gives the analytic identity
\[
 \frac{H(\epsilon^2,x_0+\epsilon y)}{u_0\epsilon^m}
   =P_m(y)+c\epsilon P_{m+1}(y)
        +d\epsilon^2P_{m+2}(y)+O(\epsilon^3).
 \tag{RJ13}
\]
The remainder is holomorphic in $(\epsilon,y)$ near every bounded
$y$-set; division by $\epsilon^m$ is legitimate because the lower
coefficients vanished. This uses the given joint analyticity and
does not assume convergence of an unrestricted heat series.

At a root $r$ the recurrences give
$P_{m+1}(r)=-2P_m'(r)$,
$P_{m+1}'(r)=0$,
$P_{m+2}(r)=-2rP_m'(r)$, and
$P_m''(r)/P_m'(r)=r/2$.
The implicit function theorem applied to RJ13 and these four
identities gives the first two root corrections $2c$ and
$(2d-c^2)r$. The complete local cluster consequently has
\[
 \theta_\nu(\epsilon)=r_\nu\epsilon+2c\epsilon^2
                   +b r_\nu\epsilon^3+O(\epsilon^4).
 \tag{RJ14}
\]
To verify completeness, take a small fixed disk in $x$ containing
only the zero at $0$ of the initial germ. On its boundary $H_0$
is nonzero, and uniform continuity and the argument principle keep
the number of zeros equal to $m$ for small $t$. The $m$ distinct
branches just constructed exhaust them. Changing $\epsilon$ to
$-\epsilon$ leaves their unordered set unchanged. Thus every
symmetric holomorphic expression in them is analytic in the original
$t$. In particular, with $s_a(t)=\sum_\nu\theta_\nu^a$,
\begin{align}
 s_{2j}(t)&=t^j\left[p_{2j}+t\left(2jb p_{2j}
                 +4j(2j-1)c^2p_{2j-2}\right)+O(t^2)\right],
                       &&j\ge1,\tag{RJ15}\\
 s_{2j+1}(t)&=2(2j+1)c p_{2j}t^{j+1}+O(t^{j+2}),
                       &&j\ge0,\tag{RJ16}\\
 s_{2j+2}(t)&=p_{2j+2}t^{j+1}+O(t^{j+2}),
                       &&j\ge0.\tag{RJ17}
\end{align}
The coefficient at degree $2j+2$ in RJ15 comes from either one
$br_\nu\epsilon^3$ or two $2c\epsilon^2$ factors, giving the
two displayed terms. The analogous single $2c\epsilon^2$ choice
gives RJ16. All odd root moments vanish by $P_m(-y)=(-1)^mP_m(y)$.
The next possible odd error powers disappear because the sums are
even in $\epsilon$, proving the stated errors in $t$. Also $s_0=m$
exactly. Newton's identity gives $p_2=2m(m-1)$, and $p_{2j}>0$
for $m\ge2$, $j\ge1$, since at least one real root is nonzero.

\paragraph{RJ6. The residue formula for every fixed-test layer.}
Let $F,G$ be fixed holomorphic functions near $\rho$, with
$f_a=F^{(a)}(\rho)/a!$, $g_a=G^{(a)}(\rho)/a!$, vanishing for
$a<j$, where $0\le j<m$. Reflection is
$F^\#(s)=\overline{F(1-\bar s)}$. In the coordinate RJ1, the
pullback coefficients are $(i/2)^a f_a$ and $(i/2)^a g_a$;
reflection conjugates their coefficients. Consequently the
coefficients of their product at $x^{2j},x^{2j+1},x^{2j+2}$ are
\[
 \frac{\bar f_jg_j}{4^j},\quad
 \frac{i(\bar f_jg_{j+1}-\bar f_{j+1}g_j)}{2\,4^j},\quad
 \frac{\bar f_{j+1}g_{j+1}-\bar f_jg_{j+2}-\bar f_{j+2}g_j}{4^{j+1}}.
 \tag{RJ18}
\]
Insert RJ15--17 into the finite branch sum
$Z_{\rho,t}(F^\#G)=\sum_\nu(F^\#G)(\rho+i\theta_\nu/2)$.
The convergent remaining Taylor series is $O(\epsilon^{2j+3})$
on each branch, and its symmetric sum is even in $\epsilon$;
it is therefore $O(t^{j+2})$. We have proved
\[
 Z_{\rho,t}(F^\#G)=\frac{t^j}{4^j}
       \left[p_{2j}\bar f_jg_j+t\mathcal C_j(F,G)+O(t^2)\right],
 \tag{RJ19}
\]
with the following residue expression for the entire HR4 coefficient:
\begin{align}
 \mathcal C_j(F,G)={}&
 \left[\left(2j p_{2j}+4j(2j-1)p_{2j-2}\right)A^2
                         -4j p_{2j}B\right]\bar f_jg_j\notag\\
 &-i(2j+1)A p_{2j}
            (\bar f_jg_{j+1}-\bar f_{j+1}g_j)\notag\\
 &+\frac{p_{2j+2}}4
       (\bar f_{j+1}g_{j+1}-\bar f_jg_{j+2}-\bar f_{j+2}g_j).
 \tag{RJ20}
\end{align}
For $j=0$ the entire first line is defined to be zero, with no
negative-index moment. Only $A=L_1/L_0$ is then needed.
For all $j\ge1$, order $3$ supplies both $A,B$. The test jets
through degree $j+2$ in RJ20 are the actual fixed-test jets. They
are not erased even if $j+2\ge m$. The formula uses the full
analytic unit through the exact coefficients it determines at this
time order; it makes no claim that a finite receiver replaces $u$
for higher orders.

\paragraph{RJ7. The residue formula for the fixed Cauchy determinant.}
Choose $m$ distinct poles $z_l$ with $\operatorname{Re}z_l>1$,
put $a_l=\rho-z_l$, $F_l(s)=1/(s-z_l)$, and let
$C_\rho(t)_{al}=Z_{\rho,t}(F_a^\#F_l)$.
With $M=m(m-1)/2$ the exact determinant is
\[
 \det C_\rho(t)=\operatorname{disc}W(t,x)\,
 \frac{4^{-M}|\prod_{a<l}(z_l-z_a)|^2}
      {|\prod_{\nu,l}(a_l+i\theta_\nu/2)|^2}
 \quad\text{for real }t.
 \tag{RJ21}
\]
Here $W=\prod_\nu(x-\theta_\nu)$ is the original monic local
factor, and its original analytic unit is retained in $H=UW$.
For a proof, the trace matrix in the basis $1,x,\ldots,x^{m-1}$
has entries $s_{a+b}$, so its determinant is the squared root
Vandermonde. The coefficient columns of the Cauchy tests are the
inverse denominators $(a_l I+(i/2)C_x)^{-1}e_0$. Evaluating them
on the distinct branches gives the Cauchy matrix. Its determinant,
after multiplication by its common denominator, is alternating in
the row points and the poles; its degrees force the two Vandermonde
factors. Induction on a successive simple-pole residue gives the
constant $(-1)^M$. Dividing by the root evaluation Vandermonde
therefore gives the coefficient determinant
\[
 \det R(t)=\frac{(-i/2)^M\prod_{a<l}(z_l-z_a)}
                     {\prod_{\nu,l}(a_l+i\theta_\nu/2)}.
\]
For real $t$, coefficient reflection gives $C_\rho=R^*KR$,
including negative $t$; it does not replace reflection by the
Euclidean evaluation norm at nonreal roots. This proves RJ21
away from zero. Analytic continuation proves it at zero as well.

The squared Vandermonde in RJ14 has first relative correction
$2Mb=m(m-1)b$. Its leading constant is
\[
 D_m=\prod_{a<l}(r_l-r_a)^2=2^M\prod_{k=1}^m k^k.
 \tag{RJ22}
\]
For completeness, $D_m=(-1)^M\prod_iP_m'(r_i)$ and
$P_m'=mP_{m-1}$. The root-product identity for monic polynomials
gives $\prod_iP_{m-1}(r_i)=\prod_jP_m(s_j)$, since $m(m-1)$
is even, where $s_j$ are the roots of $P_{m-1}$.
The recurrence makes the right side
$(-2(m-1))^{m-1}\prod_jP_{m-2}(s_j)$.
Substitution in the two discriminant formulas gives
$D_m=2^{m-1}m^mD_{m-1}$; the signs cancel because
$m-1+M-(m-1)(m-2)/2=2(m-1)$. Starting at $D_1=1$
proves RJ22.

On a local logarithm at each nonzero $a_l$, RJ16 and $p_2$
give
\[
 \sum_\nu\log(a_l+i\theta_\nu/2)
 =m\log a_l+t\left(\frac{imc}{a_l}
                    +\frac{m(m-1)}{4a_l^2}\right)+O(t^2).
\]
The real part and its derivative are independent of the logarithm
branch. We obtain the full HR13--14 expansion
\[
 \det C_\rho(t)=
 \frac{D_m4^{-M}|\prod_{a<l}(z_l-z_a)|^2}
      {\prod_l|a_l|^{2m}}
 t^M\left(1+\mathcal E_\rho t+O(t^2)\right),
 \tag{RJ23}
\]
with the exact residue correction
\[
 \boxed{\mathcal E_\rho=m(m-1)(A^2-2B)
       +2mA\operatorname{Re}\left(i\sum_l\frac1{a_l}\right)
       -\frac{m(m-1)}2\operatorname{Re}\sum_l\frac1{a_l^2}.}
 \tag{RJ24}
\]
For $m=1$ the terms with $m(m-1)$ vanish, so only $A$ occurs.
RJ24 preserves the pole configuration, every affine factor, and
the original heat time.

\paragraph{RJ8. Genuine heat counterexamples prove every threshold.}
For every $m\ge1$ and real $c,d$, define the finite polynomial
\[
 \mathsf H_{m;c,d}(t,x)=
 \sum_{a=0}^{\lfloor(m+2)/2\rfloor}
   \frac{(-t)^a}{a!}\partial_x^{2a}
                  \bigl[x^m(1+cx+dx^2)\bigr].
 \tag{RJ25}
\]
This is jointly entire, real for real arguments, and its initial
zero at $x=0$ has order exactly $m$. Termwise differentiation
and reindexing prove $\partial_t\mathsf H=-\partial_x^2\mathsf H$:
the final unpaired derivative is zero because it exceeds the
polynomial degree. Thus these are genuine analytic solutions in
the actual clock, and their unit is exactly $1+cx+dx^2$.
The cluster completeness proof in RJ5 applies; no extra initial
zeros of this polynomial are inserted into the multiplicity-$m$
cluster. These examples test information sufficiency for analytic
heat families. They are not asserted to be members of the
Rodgers--Tao family.

First vary $c$ with $d$ fixed. The order-$1$ receiver is always
$\lambda_1(1)=1$. Taking the fixed tests $F(s)=1$ and
$G(s)=s-\rho$ in RJ20 gives $\mathcal C_0(F,G)=imc$.
It changes with $c$, so order $1$ does not determine the entire
first fixed-test coefficient. Order $2$ does by RJ20. The
observation is Hermitian as well: take $F=G=1+i(s-\rho)$.
Then $\mathcal C_0(F,F)=-2mc+p_2/4$, which also changes.

Next fix $c$ and compare $d$ with $d+\Delta$, $\Delta\ne0$.
The whole order-$2$ receiver is the same in both families:
$\lambda_2(1)=-c$, $\lambda_2(x)=1$ and
$B_2=\left(\begin{smallmatrix}-c&1\\1&0\end{smallmatrix}\right)$.
For every $1\le j<m$, with the same fixed tests in both families,
RJ20 gives
\[
 \boxed{\Delta\mathcal C_j(F,G)=
              4j p_{2j}\Delta\,\bar f_jg_j,\qquad
       \Delta\mathcal E_\rho=2m(m-1)\Delta.}
 \tag{RJ26}
\]
Choose $F=G=(s-\rho)^j$ to make the first change nonzero.
This works for every $m\ge2$ and every $1\le j<m$ because
$p_{2j}>0$. The determinant change is nonzero for every allowed
pole configuration when $m\ge2$. The universal order-$3$
sufficiency in RJ20 and RJ24 is therefore sharp.

These witnesses can also be taken inside a fixed Cauchy-test
span. For the same poles define $D(s)=\prod_l(s-z_l)$ and
\[
 B_j(s)=\frac{(s-\rho)^j
                [D(\rho+v)]_{<m-j}|_{v=s-\rho}}{D(s)}.
 \tag{RJ27}
\]
The difference between its numerator and $(s-\rho)^jD(s)$
vanishes to order $m$, while $D(\rho)\ne0$. Hence
$B_j(s)=(s-\rho)^j+O((s-\rho)^m)$ and its leading jet is $1$.
Its numerator has degree at most $m-1$, so subtracting its simple
pole residues expresses it as a linear combination of the fixed
$1/(s-z_l)$; the remaining rational function has no poles and
vanishes at infinity, hence is zero. Its higher jets are fixed
by these poles and stay unchanged when $d$ varies. Thus RJ26
with $F=G=B_j$ remains $4j p_{2j}\Delta$ without dropping any
of those jets. Every flag threshold is attained by the original
fixed rational tests as well.

For $m=1$, the determinant is the only layer-$0$ one-test Gram
entry. Its correction equals
$-2c\operatorname{Re}(i/a_1)$. It needs order $2$ precisely
when $\operatorname{Re}(i/a_1)\ne0$, as varying $c$ in RJ25
proves. Since $\operatorname{Re}(i/a_1)=
\operatorname{Im}(a_1)/|a_1|^2$, it is identically zero at this
time order when $\operatorname{Im}z_1=\operatorname{Im}\rho$;
then no unit derivative is required for this particular scalar
coefficient. This exception does not change the order-$2$
threshold for the full layer-$0$ bilinear observation.

\paragraph{RJ9. What the extra datum changes at a double zero.}
For $m=2$ the original $A_2$ already determines $c=-A$.
The first larger receiver has the actual quotient and kernel
\[
 0\longrightarrow\mathbb Cx^2\longrightarrow A_3
        \xrightarrow{\pi_{3,2}}A_2\longrightarrow0,
 \qquad\lambda_2(f)=\lambda_3(x\widetilde f).
 \tag{RJ28}
\]
The new value is $L_2=\lambda_3(1)$, while
$L_1=\lambda_3(x)$ and $L_0=\lambda_3(x^2)$ were already
recorded by $A_2$. Since $p_0=2$, $p_2=4$, $p_4=8$,
the complete layer-$1$ coefficient is
\begin{align}
 \mathcal C_1(F,G)={}&16(A^2-B)\bar f_1g_1
       -12iA(\bar f_1g_2-\bar f_2g_1)\notag\\
 &+2(\bar f_2g_2-\bar f_1g_3-\bar f_3g_1),
 \tag{RJ29}\\
 \mathcal E_\rho={}&2A^2-4B
       +4A\operatorname{Re}\left(i\sum_{l=1}^2\frac1{a_l}\right)
       -\operatorname{Re}\sum_{l=1}^2\frac1{a_l^2}.
 \tag{RJ30}
\end{align}
Thus changing the new normalized datum by $\Delta B$ with the
old receiver fixed changes these coefficients by
$-16\Delta B\,\bar f_1g_1$ and $-4\Delta B$, respectively.
The actual $t^2$ coefficient in $Z_{\rho,t}$ changes by
$-4\Delta B\,\bar f_1g_1$, because RJ19 retains its factor
$1/4$. The earlier energy paper already proved the insufficiency
for its energy finite part. Equations RJ29--30 are its additional
fixed-test and determinant receivers, with the full higher test
jets included.

For an actual simple zero, $A_1$ lacks the coefficient $c$;
$A_2$ supplies it for the full first variation. For $m\ge3$ the
original $A_m$ already contains the three required residue values,
so every displayed fixed-test correction is read directly from
that original receiver. Higher values of $n$ are still related by
RJ7--9. None of these statements changes $m$ or turns a local
correction into a sign of the complete arithmetic Weil form.

\begin{figure}[p]
\centering
\begin{tikzpicture}[>=Stealth,
 box/.style={draw,rounded corners,align=center,text width=11.4cm,
 inner sep=8pt}]
\node[box] (three) at (0,0)
 {$A_3$, marked $x$, full $\lambda_3$\\
 $q_0=\lambda_3(x^2)$, $q_1=\lambda_3(x)$,
 $q_2=\lambda_3(1)$};
\node[box] (two) at (0,-2.2)
 {$A_2$, marked $x$, full $\lambda_2$\\
 $q_0=\lambda_2(x)$, $q_1=\lambda_2(1)$};
\node[box] (one) at (0,-4.4)
 {$A_1$, full $\lambda_1$\\$q_0=\lambda_1(1)=u_0^{-1}$};
\draw[->] (three)--node[right,align=left]
 {$\pi_{3,2}$; kernel $\mathbb Cx^2$}(two);
\draw[->] (two)--node[right,align=left]
 {$\pi_{2,1}$; kernel $\mathbb Cx$}(one);
\node[box] (receiver) at (0,-7.2)
 {$A=q_1/q_0$, $B=q_2/q_0$\\
 $c=-A$, $d=A^2-B$, $b=A^2-2B$\\
 Layer $j=0$: $A$; every $j\ge1$: $A,B$\\
 Determinant for $m\ge2$: $A,B$};
\end{tikzpicture}
\caption{The coordinate-marked residue tower and the sharp fixed-test
receivers. Each quotient's transpose under the perfect residue forms
is multiplication by $x$, not an algebra inclusion. RJ3 proves these
maps and their kernels; RJ6--9 prove all displayed receiver thresholds
using the original $t$ and $s=\rho+ix/2$. The energy receiver is credited
to the supplied Theorem \texttt{thm:residuejet}; the fixed-test targets
are HR4 and HR14. The diagram records analytic thickenings and does not
assign a zero multiplicity.}
\end{figure}
